\section*{Hoofdstuk \ref{ch:b3:exoplanets} --- Exoplaneten}

\begin{solution}{exo:b3:exoplanets:1}
(a)~$K = (2\pi G/P)^{1/3}m_p/\bigl((M_\odot + m_p)^{2/3}\sqrt{1 - e^2}\bigr)$.
Jupiter: $P = \qty{3.742e8}{s}$, $(2\pi G/P)^{1/3} = \num{1.039e-6}$ in SI-eenheden
, $K = \num{1.039e-6}\times\num{1.898e27}/(\num{1.990e30})^{2/3}/0.9988
= \qty{12.5}{m.s^{-1}}$. De Aarde: $P = \qty{3.155e7}{s}$, $K =
\qty{0.089}{m.s^{-1}}$, ongeveer \qty{9}{cm.s^{-1}}. (b)~\qty{3.5}{km/h} $=
\qty{0.97}{m.s^{-1}}$ en \qty{0.35}{km/h} $= \qty{0.097}{m.s^{-1}}$.
Het signaal van Jupiter is dertien keer de eerste precisie; dat van de Aarde
ligt zelfs onder de tweede, en de eigen ruis van de ster komt er bovenop.
\end{solution}

\begin{solution}{exo:b3:exoplanets:2}
(a)~$(71492/695700)^2 = \num{0.0106}$, een helling van \qty{1.06}{\percent};
$(6378/695700)^2 = \num{8.4e-5}$, 84 delen per miljoen. (b)~$R_\star/a$:
de Aarde $695700/\num{1.496e8} = \num{4.7e-3}$, ongeveer één kans op 215;
Jupiter $695700/\num{7.785e8} = \num{8.9e-4}$, één kans op 1100.
\end{solution}

\begin{solution}{exo:b3:exoplanets:3}
$m_p\sin i \approx K M_\odot^{2/3}(P/2\pi G)^{1/3}$ met $P =
\qty{3.655e5}{s}$: $(P/2\pi G)^{1/3} = \num{9.55e4}$ in SI-eenheden, dus $m_p
\sin i = 59\times(\num{1.9884e30})^{2/3}\times\num{9.55e4} =
\qty{8.9e26}{kg}$, dat is 0.47 Jupitermassa's of 149 aardmassa's (de
correctie voor $m_p$ in $M_\star + m_p$ ligt onder één deel op duizend).
$a = \bigl(GM_\odot P^2/4\pi^2\bigr)^{1/3} = \qty{7.66e9}{m} =
\qty{0.051}{au}$.
\end{solution}

\begin{solution}{exo:b3:exoplanets:4}
(a)~$4GM/c^2 = 4\times\num{6.674e-11}\times\num{9.94e29}/(\num{2.998e8})^2 =
\qty{2953}{m}$; $(D_S - D_L)/(D_LD_S) = 1/D_S\times(D_S - D_L)/D_L =
1/(\qty{8}{kpc}) = 1/\qty{2.469e20}{m}$. $\theta_\mathrm E =
\sqrt{2953/\num{2.469e20}} = \qty{3.46e-9}{rad} = \qty{0.71}{mas}$. Bij de
lens, $R_\mathrm E = \theta_\mathrm E D_L = \num{3.46e-9}\times
\qty{1.234e20}{m} = \qty{4.27e11}{m} = \qty{2.85}{au}$. (b)~$A(1) =
3/\sqrt5 = 1.34$; $A(0.5) = 2.25/(0.5\times2.062) = 2.18$; $A(0.1) =
2.01/(0.1\times2.002) = 10.0$. (c)~$t_\mathrm E = \num{4.27e11}/\num{2e5}
= \qty{2.13e6}{s}$, ongeveer 25 dagen: microlensing-gebeurtenissen duren weken.
\end{solution}

\begin{solution}{exo:b3:exoplanets:5}
(a)~$a = (GM_\odot P^2/4\pi^2)^{1/3} = \qty{6.09e9}{m}$ (\qty{0.0407}{au}),
dus $a/R_\star = 8.75$. (b)~Met $P/\pi = \qty{22.9}{h}$, voor $b = 0$:
\[
T_{14} = \frac P\pi\arcsin\frac{1.12}{8.75} = \qty{2.94}{h}, \qquad
T_{23} = \frac P\pi\arcsin\frac{0.88}{8.75} = \qty{2.31}{h}.
\]
Voor $b = 0.7$: $\cos i = 0.7/8.75$, $i = \ang{85.4}$, $(a/R_\star)\sin i =
8.726$, en
\[
T_{14} = \frac P\pi\arcsin\frac{\sqrt{1.2544 - 0.49}}{8.726} = \qty{2.30}{h},
\qquad
T_{23} = \frac P\pi\arcsin\frac{\sqrt{0.7744 - 0.49}}{8.726} = \qty{1.40}{h}.
\]
(c)~De intrede en uittrede nemen $(T_{14} - T_{23})/T_{14} = \qty{21}{\percent}$
van de transit wanneer die centraal is, \qty{39}{\percent} bij $b = 0.7$; een
randtransit, $b \to 1$, heeft helemaal geen vlakke bodem en ziet eruit als een V.
\end{solution}

\begin{solution}{exo:b3:exoplanets:6}
$\bar\rho = M/\tfrac43\pi R^3$, met $\tfrac43\pi R_\oplus^3 =
\qty{1.087e21}{m^3}$: TRAPPIST-1~b $1.374\times\num{5.97e24}/(1.116^3\times
\num{1.087e21}) = \qty{5.43}{g.cm^{-3}}$; Kepler-10~b
\qty{5.60}{g.cm^{-3}}; WASP-39~b \qty{0.167}{g.cm^{-3}}; HD~209458~b
\qty{0.337}{g.cm^{-3}}. De eerste twee zijn even dicht als de Aarde
(\qty{5.49}{g.cm^{-3}} met dezelfde constanten) en rotsachtig; de laatste twee
zijn gasreuzen, lichter dan water, WASP-39~b acht keer lichter dan
Jupiter.
\end{solution}

\begin{solution}{exo:b3:exoplanets:7}
De extremen zijn $\gamma + K(1 + e\cos\omega)$ en $\gamma - K(1 -
e\cos\omega)$. Hun verschil is $2K = \qty{60.0}{m.s^{-1}}$, dus $K =
\qty{30.0}{m.s^{-1}}$; hun gemiddelde is $\gamma + Ke\cos\omega =
\qty{20.4}{m.s^{-1}}$, dus $Ke\cos\omega = \qty{8.0}{m.s^{-1}}$ en
$e\cos\omega = 0.27$. Bij een excentrische baan vertoeft de ster lang nabij
het apastron en snelt door het periastron, dus de kromme is scheef en het
middelpunt van de extremen is niet het tijdgemiddelde $\gamma$.
\end{solution}

\begin{solution}{exo:b3:exoplanets:8}
(a)~$g = GM/R^2 = \num{6.674e-11}\times\num{5.332e26}/(\num{9.144e7})^2 =
\qty{4.26}{m.s^{-2}}$; $H = kT/\bar m g = \num{1.381e-23}\times1100/(2.3
\times\num{1.661e-27}\times4.26) = \qty{934}{km}$. Met $R_\star =
\qty{6.533e8}{m}$, de diepte is $(\num{9.144e7}/\num{6.533e8})^2 =
\num{0.0196}$. (b)~$2HR_p/R_\star^2 = 2\times\num{9.34e5}\times
\num{9.144e7}/(\num{6.533e8})^2 = \num{4.0e-4}$, 400 delen per miljoen;
twee schaalhoogten, 800: een gasreus met lage zwaartekracht is het gemakkelijkste doelwit
voor \omterm{def:b3:exoplanets:transmission}{transmissiespectroscopie}.
\end{solution}

\begin{solution}{exo:b3:exoplanets:9}
$p = 1.116\times6378/(0.1192\times695700) = 0.0858$, $p^2 = \num{7.37e-3}$.
$B_\lambda \propto 1/(e^{hc/\lambda kT} - 1)$ met $hc/\lambda k =
\qty{959}{K}$ bij \qty{15}{\micro m}: de ster $1/(e^{0.374} - 1) = 2.20$;
een dagzijde bij \qty{400}{K}, $1/(e^{2.40} - 1) = 0.0998$, verhouding \num{0.0453},
diepte \num{3.3e-4} (334 delen per miljoen); bij \qty{500}{K}, $1/(e^{1.92} -
1) = 0.172$, verhouding \num{0.078}, diepte \num{5.7e-4}. Honderd kelvin
verandert de diepte met \qty{70}{\percent}: de gemeten diepte legt de
temperatuur van de dagzijde nauwkeurig vast, en een hete dagzijde betekent dat weinig warmte
naar de nachtzijde wordt getransporteerd.
\end{solution}

\begin{solution}{exo:b3:exoplanets:10}
(a)~$S \propto L/d^2$, dus $d = \sqrt{0.005}\,\unit{au} = \qty{0.0707}{au}$;
$P = 2\pi\sqrt{a^3/GM} = \qty{1.33e6}{s} = \qty{15.4}{d}$.
(b)~$K = (2\pi G/P)^{1/3}M_\oplus/(0.2M_\odot)^{2/3} = \qty{0.75}{m.s^{-1}}$,
8.4 keer het effect van de Aarde op de Zon: de lichte ster beweegt meer, en
de korte periode helpt ook mee. (c)~$(6378/139140)^2 = \num{2.1e-3}$, 25 keer
de transit van de Aarde voor de Zon. Beide signalen zijn groter, de transits
komen elke twee weken, en de transitkans $R_\star/a =
\num{0.013}$ is drie keer die van de Aarde: kleine sterren geven rotsplaneten in
hun bewoonbare zones een kans om gevonden en bestudeerd te worden.
\end{solution}

\begin{solution}{exo:b3:exoplanets:11}
Met $r\,dr = -\mu\,d\mu$ is de gemiddelde intensiteit $\frac1\pi\int_0^1
I\,2\pi r\,dr = 2\int_0^1 I(\mu)\mu\,d\mu = 2\bigl[\tfrac12 - u_1(\tfrac12 -
\tfrac13) - u_2\int_0^1(1 - \mu)^2\mu\,d\mu\bigr]$. De laatste integraal is
$\tfrac12 - \tfrac23 + \tfrac14 = \tfrac1{12}$, dus het gemiddelde is $1 - u_1/3 -
u_2/6$, hier $0.825$. Een kleine planeet verbergt $\pi p^2 I(\mu)$ van $\pi\times
0.825$: diepte $p^2 I(\mu)/0.825$. Bij $r = 0$, $I = 1$: $1.212\,p^2$; bij
$r = 0.5$ ($\mu = 0.866$), $I = 0.942$: $1.142\,p^2$; bij $r = 0.9$ ($\mu =
0.436$), $I = 0.695$: $0.842\,p^2$. Een uniforme-schijfmodel zou
$1.212\,p^2$ lezen als $p^2$ en de straal overschatten met $\sqrt{1.212} = 1.10$,
\qty{10}{\percent}.
\end{solution}

\begin{solution}{exo:b3:exoplanets:12}
(a)~$R_\odot/a$: \num{9.3e-3} bij \qty{0.5}{au}, \num{6.6e-3} bij
\qty{0.7}{au}, \num{4.7e-3} bij \qty{1}{au}. (b)~$42000\times0.11\times
\num{6.6e-3} = 31$ planeten transiteren. (c)~Slechts ongeveer een derde ervan werd
gevonden: de overige gaven transits die te ondiep waren (84 delen per miljoen voor
een Aarde) of te weinig (één of twee per jaar) om boven de ruis uit te komen. Die
fractie is de detectie-efficiëntie; zonder daarvoor te corrigeren, en voor de
transitkans, zouden de tellingen de frequentie van dergelijke planeten onderschatten
met een factor van ongeveer 460.
\end{solution}

\begin{solution}{pb:b3:exoplanets:1}
\textbf{1.} $M_\star = 0.80\times\num{1.9884e30} = \qty{1.591e30}{kg}$, $P =
\qty{1.0368e6}{s}$: $a = (GM_\star P^2/4\pi^2)^{1/3} = \qty{1.4245e10}{m}
= \qty{0.0952}{au}$.

\textbf{2.} $v = 2\pi a/P = \qty{86.3}{km.s^{-1}}$.

\textbf{3.} $m_p\sin i = KM_\star^{2/3}(P/2\pi G)^{1/3} = 2.59\times\num{1.363e20}
\times\num{1.352e5} = \qty{4.77e25}{kg}$, \num{7.99} aardmassa's.

\textbf{4.} $a_\star = a\,m_p/M_\star = \num{1.4245e10}\times\num{4.77e25}/
\num{1.591e30} = \qty{427}{km}$; $2\pi a_\star/P = \qty{2.59}{m.s^{-1}}$,
de waargenomen $K$.

\textbf{5.} $2.59$ keer een precisie van \qty{1}{m.s^{-1}}, 26 keer een van
\qty{0.1}{m.s^{-1}}: de planeet ligt binnen bereik van de beste instrumenten,
nauwelijks van die van de vorige generatie.

\textbf{6.} $m_p\sin i \propto K\sqrt{1 - e^2}$: $\times0.980$, 7.83 aardmassa's,
\qty{2}{\percent} minder.

\textbf{7.} $p = \sqrt{\num{8.38e-4}} = 0.02895$.

\textbf{8.} $R_\star = 0.76\times\qty{695700}{km} = \qty{528700}{km}$; $R_p =
\qty{15310}{km} = 2.40$ aardstralen.

\textbf{9.} $a/R_\star = \num{1.4245e10}/\num{5.287e8} = 26.94$.

\textbf{10.} $T_{14} = (\qty{288}{h}/\pi)\arcsin(1.02895/26.94) = \qty{3.50}{h}$.

\textbf{11.} $\sqrt{(1 + p)^2 - b^2} = (a/R_\star)\sin(\pi T_{14}/P) = 26.94\times
\sin(\pi\times3.23/288) = 0.949$ (met $\sin i \approx 1$); $b^2 =
1.0587 - 0.9008 = 0.158$, $b = 0.40$.

\textbf{12.} $\cos i = b\,R_\star/a = 0.40/26.94 = 0.0148$, $i = \ang{89.15}$.

\textbf{13.} $R_\star/a = 1/26.94 = 0.037$: één kans op 27.

\textbf{14.} $\sin i = 0.99989$: de echte massa is 7.99 aardmassa's, de \omterm{def:b3:exoplanets:rv}{minimale
massa} binnen \num{1e-4}.

\textbf{15.} $\bar\rho = \num{4.77e25}/\bigl(\tfrac43\pi(\num{1.531e7})^3\bigr) =
\num{4.77e25}/\num{1.502e22} = \qty{3180}{kg.m^{-3}}$,
\qty{3.18}{g.cm^{-3}}.

\textbf{16.} De planeet is $2.40/1.74 = 1.38$ keer groter dan een rotsplaneet van
dezelfde massa, 2.6 keer het volume. Gesteente alleen kan dat niet verklaren: er is
een omhulling van licht gas nodig, enkele procenten van de massa in waterstof en
helium, of een groot aandeel water.

\textbf{17.} Bij 2.40 aardstralen ligt hij boven het dal, in de piek van 2--3 aardstralen:
een \omterm{def:b3:exoplanets:classes}{sub-Neptunus}.

\textbf{18.} $g = Gm_p/R_p^2 = \num{6.674e-11}\times\num{4.77e25}/(\num{1.531e7})^2 =
\qty{13.6}{m.s^{-2}}$.

\textbf{19.} $L = 4\pi R_\star^2\sigma T^4 = 4\pi(\num{5.287e8})^2\times\num{5.670e-8}
\times5000^4 = \qty{1.245e26}{W}$, 0.325 zonneluminositeiten.

\textbf{20.} $S/S_\oplus = 0.325/0.0952^2 = 35.9$: ver binnen de binnenrand van de
bewoonbare zone.

\textbf{21.} $L/4\pi a^2 = \qty{4.88e4}{W.m^{-2}}$; $T_\mathrm{eq} =
\bigl(\num{4.88e4}\times0.7/(4\sigma)\bigr)^{1/4} = \qty{623}{K}$.

\textbf{22.} $H = kT/\bar m g = \num{1.381e-23}\times623/(2.3\times\num{1.661e-27}
\times13.6) = \qty{166}{km}$.

\textbf{23.} $2HR_p/R_\star^2 = 2\times\num{1.66e5}\times\num{1.531e7}/
(\num{5.287e8})^2 = \num{1.82e-5}$, 18 delen per miljoen; een band van twee
schaalhoogten hoog, 36 delen per miljoen.

\textbf{24.} $H$ is $18/2.3 = 7.8$ keer kleiner, \qty{21}{km}: 2.3 delen per
miljoen per schaalhoogte, 4.6 voor de band. Bij tien delen per miljoen toont de
waterstofrijke atmosfeer haar waterband op meer dan drie keer de
ruis; een stoomatmosfeer ziet er vlak uit, net als een bewolkte of een planeet
zonder lucht.

\textbf{25.} \textbf{\qty[text-series-to-math=true]{3.2}{g/cm^3}}.
\end{solution}
